\section{Pattern 3：Reviewer（审查者）}

Reviewer 模式将\textbf{检查什么}与\textbf{如何检查}分离。与其在 system prompt 中写一长串代码质量规则，不如将模块化的评分标准（rubric）存储在 \texttt{references/review-checklist.md} 文件中。

\begin{importantbox}{Reviewer 的核心设计}
指令保持静态不变，Agent 动态加载外部 checklist 中的具体审查标准，并按严重级别输出结构化的评审结果。更换 checklist 即可切换审查维度：Python 风格检查 $\rightarrow$ OWASP 安全审计，无需修改 skill 基础设施。
\end{importantbox}

\subsection{工作原理}

\begin{enumerate}
    \item 用户提交代码
    \item Agent 加载 \texttt{references/review-checklist.md} 中的检查清单
    \item Agent 逐项对照清单进行系统化评审
    \item 按严重级别（severity）分组输出发现的问题
\end{enumerate}

\subsection{适用场景}

\begin{itemize}
    \item 自动化 PR review
    \item 安全漏洞预检
    \item 代码质量评估
    \item 合规性检查
\end{itemize}

\begin{knowledgebox}{模块化的威力}
Reviewer 模式最大的优势在于 checklist 的可插拔性。同一套 skill 基础设施，配合不同的 checklist 文件，可以变成完全不同的专业审查工具。这种设计让 skill 的复用效率最大化。
\end{knowledgebox}

\subsection{本章小结}

Reviewer 通过将审查标准外置为独立的 checklist 文件，实现了检查逻辑与检查内容的解耦。更换 checklist 即可获得全新的审查能力，是自动化代码审查和安全审计的理想选择。

